% Part: normal-modal-logic
% Chapter: syntax-and-semantics
% Section: introduction

\documentclass[../../../include/open-logic-section]{subfiles}

\begin{document}

\olfileid{nml}{syn}{int}

\olsection{مقدمه}

منطق موجهات به \emph{گزاره‌های موجهاتی} و روابط استلزام میان آن‌ها
می‌پردازد. نمونه‌هایی از گزاره‌های موجهاتی عبارت‌اند از:
\begin{enumerate}
\item ضرورتاً~$2+2=4$.
\item ضرورتاً ممکن است فردا باران ببارد.
\item اگر ضرورتاً ممکن باشد که~$!A$، آنگاه ممکن است که~$!A$.
\end{enumerate}
امکان و ضرورت تنها وجه‌ها نیستند؛ رابط‌های یکانی دیگری نیز وجه شمرده
می‌شوند، برای نمونه «باید چنین باشد که~$!A$»، «چنین خواهد شد که~$!A$»،
«دانا می‌داند که~$!A$» یا «دانا باور دارد که~$!A$».

منطق موجهات نخستین بار در \emph{دربارهٔ عبارت} ارسطو ظاهر می‌شود. او
نخستین کسی بود که دریافت ضرورت مستلزم امکان است، اما عکس آن برقرار
نیست؛ امکان و ضرورت برحسب یکدیگر تعریف‌پذیرند؛ اگر
$!A \land !B$ ممکن‌الصدق باشد، آنگاه~$!A$ و~$!B$ هر یک ممکن‌الصدق‌اند،
اما عکس آن برقرار نیست؛ و اگر~$!A \to !B$ ضروری باشد، آنگاه اگر~$!A$
ضروری باشد، $!B$ نیز ضروری است.

نخستین رویکرد مدرن به منطق موجهات کار سی.~آی. لوئیس بود که با کتاب
لوئیس و لنگفورد، \emph{منطق نمادی} (۱۹۳۲)، به اوج رسید. لوئیس و
لنگفورد از نمایش استلزام به‌وسیلهٔ شرطی مادی خشنود نبودند:
$!A \lif !B$ جایگزین نامناسبی برای «$!A$ مستلزم~$!B$ است» به شمار
می‌رفت. آن‌ها در عوض پیشنهاد کردند استلزام را به‌صورت «ضرورتاً، اگر
$!A$ آنگاه~$!B$» مشخص کنیم، که با~$!A \strictif !B$ نمادگذاری می‌شود.
لوئیس در تلاش برای تفکیک خواص گوناگون، پنج دستگاه موجهاتی متفاوت،
\Log{S1}، \ldots، \Log{S4} و~\Log{S5} را شناسایی کرد که دو دستگاه آخر
هنوز به کار می‌روند.

رویکرد لوئیس و لنگفورد کاملاً \emph{نحوی} بود: آن‌ها اصل‌ها و قواعد
معقولی را شناسایی و بررسی کردند با این ابزارها چه چیزهایی اثبات‌پذیر
است. رویکرد معناشناختی مدت‌ها دور از دسترس ماند تا اینکه رودلف کارناپ
در \emph{معنا و ضرورت} (۱۹۴۷) نخستین تلاش را با استفاده از مفهوم
\emph{توصیف حالت} انجام داد؛ یعنی مجموعه‌ای از جمله‌های اتمی، همان
جمله‌هایی که در آن توصیف حالت «صادق»اند. کارناپ پس از گسترش تعریف صدق
به جمله‌های دلخواه~$!A$، $!A$ را هنگامی \emph{ضرورتاً صادق} تعریف می‌کند
که در همهٔ توصیف‌های حالت صادق باشد. رویکرد کارناپ نمی‌توانست وجه‌های
\emph{تکرارشده} را مدیریت کند، زیرا جمله‌هایی به‌صورت «ممکن است ضرورتاً
\ldots ممکن باشد که~$!A$» همواره به درونی‌ترین وجه فروکاسته می‌شدند.

جهش بزرگ در معناشناسی موجهات با مقالهٔ سول کریپکی، «قضیه‌ای تمامیت در
منطق موجهات» (JSL، ۱۹۵۹)، رخ داد. کریپکی کارش را بر اندیشهٔ لایبنیتس
بنا کرد که یک گزاره ضرورتاً صادق است اگر «در همهٔ جهان‌های ممکن» صادق
باشد. بااین‌حال، این اندیشه همان کاستی رویکرد کارناپ را دارد، زیرا صدق
یک گزاره در جهان~$w$، یا توصیف حالت~$s$، اصلاً به~$w$ وابسته نیست. پس
کریپکی فرض کرد جهان‌ها با \emph{رابطهٔ دسترسی}~$R$ به یکدیگر مرتبط‌اند
و گزاره‌ای به‌صورت «ضرورتاً~$!A$» در جهان~$w$ صادق است اگر و تنها اگر
$!A$ در همهٔ جهان‌های~$w'$ که \emph{از~$w$ دسترس‌پذیرند} صادق باشد.
معناشناسی‌هایی که گونه‌ای از این رویکرد را فراهم می‌کنند معناشناسی
کریپکی نامیده می‌شوند و رشد پرشتاب منطق‌های موجهات را ممکن ساختند.

منطق موجهات، هنگامی که با معناشناسی کریپکی تفسیر شود، نشان می‌دهد
\emph{ساختارهای رابطه‌ای} «از درون» چگونه‌اند. ساختار رابطه‌ای صرفاً
مجموعه‌ای مجهز به یک رابطهٔ دوتایی است؛ برای نمونه، مجموعهٔ دانش‌آموزان
کلاس که با شمارهٔ شناسایی اجتماعی مرتب شده‌اند، ساختاری رابطه‌ای است.
اما ساختارهای رابطه‌ای در دامنه‌های بسیار گوناگونی پدید می‌آیند: افزون
بر امکان نسبی حالت‌های جهان، می‌توانیم حالت‌های معرفتی یک عامل را داشته
باشیم که با امکان معرفتی مرتبط‌اند، یا حالت‌های یک دستگاه پویا را با
گذارهای حالتشان، و مانند آن. منطق موجهات می‌تواند همهٔ این موارد را
مدل کند: مورد نخست منطق موجهات عادی یا صدقی، و موارد دیگر منطق معرفتی،
منطق پویا و جز آن را به دست می‌دهند.

بر زاویه‌ای خاص تمرکز می‌کنیم که منطقدانان موجهات آن را «نظریهٔ تناظر»
می‌نامند. یکی از مهم‌ترین کشف‌های آغازین کریپکی این است که بسیاری از
خواص رابطهٔ دسترسی~$R$، مانند تعدی یا تقارن، را می‌توان \emph{در خود
زبان موجهات} با «طرحواره‌های موجهاتی» مناسب مشخص کرد. برای نمونه،
منطقدانان موجهات می‌گویند بازتابی‌بودن~$R$ با طرحوارهٔ «اگر ضرورتاً
$!A$، آنگاه~$!A$» «تناظر» دارد. ما عمدتاً نظریهٔ تناظر چند دستگاه
کلاسیک منطق موجهات، مانند~\Log{S4} و~\Log{S5}، را بررسی می‌کنیم که با
ترکیبی از طرحواره‌های~\Ax{D}، \Ax{T}، \Ax{B}، \Ax{4} و~\Ax{5} به دست
می‌آیند.

\end{document}
